\section{Редукция гипотезы о близнецах к одному узлу}\label{sec:reduction}

Этот раздел --- стержень статьи. Мы сжимаем всю близнецовую ветвь в одну машинно-проверенную цепь
и на её дальнем конце обнажаем единственный открытый узел, на котором всё держится, --- и
единственную заражённую \code{sorry} декларацию \code{twin\_prime\_conjecture}, которая его
фиксирует. Всюду мы держим дисциплину введения: это \emph{редукция}, а не доказательство. Каждый
транзитный шаг ниже \green{} машинно проверен и свободен от распределения; вся трудность заперта в
одном \red{} утверждении, и мы нигде не выдаём редукцию за закрытие гипотезы.

\subsection{Мост к бесконечности}

Мы работаем с центровой формой пары близнецов. Всякая пара близнецов, кроме $(3,5)$, имеет вид
$(6m-1,6m+1)$, так что центровая кодировка ничего не теряет.

\begin{definition}[\code{IsTwinCenter}, \code{TwinLowers}]\label{def:istwincenter}
Для $m\in\N$
\[
  \mathrm{IsTwinCenter}(m)\;:\Longleftrightarrow\;(6m-1)\ \text{простое}\ \wedge\ (6m+1)\ \text{простое},
\]
и $\mathrm{TwinLowers}=\{\,p\mid p\ \text{и}\ p+2\ \text{простые}\,\}$. Гипотеза о близнецах --- это
в точности \code{TwinLowers.Infinite}.
\end{definition}

Мост от блочной комбинаторики к арифметике собирается из трёх элементарных нециклических звеньев.
Первое --- чистая мощность.

\begin{theorem}[\code{survivor\_of\_not\_covered}]\label{thm:survivor_of_not_covered}
Для конечных $\Omega,\mathrm{bad}\subseteq\N$
\begin{equation}\label{eq:survivor}
  |\mathrm{bad}|<|\Omega|\;\Longrightarrow\;\exists\,m\in\Omega,\ m\notin\mathrm{bad}.
\end{equation}
\end{theorem}

\emph{Идея доказательства.} Принцип Дирихле в чистейшем виде: будь $\Omega\subseteq\mathrm{bad}$,
монотонность мощности (\code{Finset.card\_le\_card}) дала бы $|\Omega|\le|\mathrm{bad}|$, вопреки
посылке. Ни анализа, ни распределения простых --- только \code{Finset} и порядок на $\N$. \green{}

Второе звено --- классический критерий простоты и его немедленное следствие для центра.

\begin{theorem}[\code{prime\_of\_no\_small\_prime\_factor}]\label{thm:prime_of_no_small_prime_factor}
Пусть $n\ge 2$ и $\forall p\ \text{простое},\ p^2\le n\Rightarrow p\nmid n$. Тогда $n$ простое.
\end{theorem}

\emph{Идея доказательства.} От противного через наименьший простой делитель $p=\mathrm{minFac}(n)$:
минимальность даёт $p\le n/p$ (\code{Nat.minFac\_le\_div}), откуда $p^2\le n$, вопреки посылке.
Наименьший простой делитель составного числа не превосходит $\sqrt n$, так что решето «до корня»
решает вопрос простоты. \green{}

\begin{theorem}[\code{isTwinCenter\_of\_root\_sieve}]\label{thm:isTwinCenter_of_root_sieve}
Если $6m\pm1\ge 2$ и ни одна сторона не делится ни на одно простое до своего корня, то
$\mathrm{IsTwinCenter}(m)$.
\end{theorem}

\emph{Идея доказательства.} Два применения \cref{thm:prime_of_no_small_prime_factor}, по одному на
сторону, упакованные в конъюнкцию \code{IsTwinCenter}. Направление \emph{выживший $\Rightarrow$
близнец} не содержит гипотез вовсе: оно элементарно и потому фиксирует важную границу --- вся
содержательная трудность лежит не здесь, а в том, чтобы \emph{произвести} выжившего, просеянного до
корня. \green{}

Третье звено --- теоретико-порядковое.

\begin{theorem}[\code{infinite\_of\_unbounded\_centers}]\label{thm:infinite_of_unbounded_centers}
Если при каждом $N$ есть центр $m>N$ с $\mathrm{IsTwinCenter}(m)$, то \code{TwinLowers.Infinite}.
\end{theorem}

\emph{Идея доказательства.} Через \code{Set.infinite\_of\_not\_bddAbove}: свидетель
$6m-1\in\mathrm{TwinLowers}$ превосходит $N$, а его сосед $6m+1$ прост, так что \code{TwinLowers}
не ограничено сверху, а значит бесконечно. Монотонность $6m-1$ по $m$ превращает «сколь угодно
высокий центр» в «сколь угодно большой нижний близнец». \green{}

\subsection{Выделение блочного ядра}

Три звена склеиваются в локальную лемму переноса и глобальную редукцию. Зафиксируем на масштабе $N$
конечный \emph{носитель} центров-кандидатов выше $N$ и конечное множество \emph{bad} центров,
заведомо не являющихся центрами близнецов.

\begin{theorem}[\code{twin\_center\_of\_block}]\label{thm:twin_center_of_block}
Пусть $\mathit{carrier},\mathit{bad}\subseteq\N$ конечны и
\begin{equation}\label{eq:block}
  (\forall m\in\mathit{carrier},\ N<m),\quad |\mathit{bad}|<|\mathit{carrier}|,\quad
  (\forall m\in\mathit{carrier},\ m\notin\mathit{bad}\Rightarrow\mathrm{IsTwinCenter}(m)).
\end{equation}
Тогда $\exists m,\ N<m\ \wedge\ \mathrm{IsTwinCenter}(m)$.
\end{theorem}

\emph{Идея доказательства.} \cref{thm:survivor_of_not_covered} извлекает выжившего
$m\in\mathit{carrier}$ с $m\notin\mathit{bad}$; первая конъюнкция ставит его выше $N$, третья делает
центром близнецов. Две строчки комбинаторики \code{Finset}, без апелляции к плотности простых.
\green{}

\begin{theorem}[\code{twin\_prime\_conjecture\_of\_blocks}]\label{thm:twin_prime_conjecture_of_blocks}
Если при каждом $N$ существуют конечные $\mathit{carrier},\mathit{bad}$, удовлетворяющие
\cref{eq:block}, то \code{TwinLowers.Infinite}.
\end{theorem}

\emph{Идея доказательства.} \cref{thm:infinite_of_unbounded_centers} сводит цель к неограниченности
центров близнецов; при каждом $N$ блочная посылка вместе с \cref{thm:twin_center_of_block} даёт
центр близнецов выше $N$. \green{} Импликация нециклична: она нигде не использует гипотезу и не
прячет её в посылку. Открытой остаётся ровно посылка --- строгое неравенство мощностей плюс критерий
выживания, то есть чисто комбинаторная задача о конечных множествах. Узел (3), выживший
$\Rightarrow$ близнец, уже закрыт \cref{thm:isTwinCenter_of_root_sieve}; остаётся \emph{счётное
ядро} --- перевес носителя над bad.

\subsection{Первый именованный узел: реальный four-corner}

Счётное ядро сперва кристаллизуется как строгий \emph{реальный four-corner} на счётах рангов.
Обозначая $R_{ij}$ множества рангов, гипотеза $H$ (\code{FourCornerInput}) требует при каждом $N$:
$0<|R_{00}|$, строгий four-corner $|R_{00}|\,|R_{33}|<|R_{03}|\,|R_{30}|$, лёгкий побочный угол
$|R_{03}|\,|R_{30}|\le|R_{00}|^2$, вместе с $|\mathit{carrier}|=|R_{00}|$,
$|\mathit{bad}|=|R_{33}|$ и носителевыми/близнецовыми условиями \cref{eq:block}. Через
\code{N33\_lt\_N00\_of\_four\_corner} они дают $|R_{33}|<|R_{00}|$, то есть в точности
$|\mathit{bad}|<|\mathit{carrier}|$.

\begin{theorem}[\code{twin\_primes\_of\_four\_corner}]\label{thm:twin_primes_of_four_corner}
\green{} (зелёная условная теорема: чистая по аксиомам, условная относительно именованного входа $H$)
\begin{equation}\label{eq:fourcorner}
  H\;\Longrightarrow\;\code{TwinLowers.Infinite}.
\end{equation}
\end{theorem}

\emph{Идея доказательства.} Развернуть $H\,N$, вывести $|\mathit{bad}|<|\mathit{carrier}|$ из
four-corner и применить \cref{thm:twin_prime_conjecture_of_blocks}. \code{sorry} нет: заключение
следует из явного входа $H$, и всё нетривиальное содержание несёт $H$. Контрапозиция фиксирует тот
же факт вручную.

\begin{theorem}[\code{twin\_finite\_contradiction}]\label{thm:twin_finite_contradiction}
Если $\neg\,\code{TwinLowers.Infinite}$ и \code{FourCornerInput} $H$, то \code{False}.
\end{theorem}

\emph{Идея доказательства.} Буквально $\mathrm{hfin}\,(\code{twin\_primes\_of\_four\_corner}\,H)$:
конечность --- отображение из бесконечности в пустоту, и мы скармливаем ей ровно ту бесконечность,
которую производит four-corner. \green{} Логически контрапозиция ничего не добавляет, но делает
осязаемым, что ломается при отрицании тезиса: конечный хвост должен быть покрыт bad-центрами без
выжившего, что $H$ запрещает.

\begin{remark}
Три независимых маршрута --- объединительная оценка по простым, бесконечный спуск EPMI и four-corner
--- были прощупаны, и все сводятся к одному объекту: строгий реальный four-corner плюс нижняя оценка
носителя. Запас $|R_{03}|\,|R_{30}|-|R_{00}|\,|R_{33}|$, нормированный масштабом, численно ползёт к
нулю: модельный four-corner (\code{model\_four\_corner}, доказан безусловно) \emph{не} переносится в
реальность свободно от распределения. Это проблема чётности, и она заперта целиком внутри \red{}
входа $H$.
\end{remark}

\subsection{Платёжный гроссбух и закон дефекта}

Остановка не бесплатна: у неё есть алгебраическая цена, доказуемая лишь на \code{Nat}/\code{Int} и
\code{Finset}. Свободный проход к упорядоченному простому $p$ требует беспошлинности у всех меньших
малых простых, так что сдвиг $a-\theta$ (при $\theta=\sigma\varepsilon$) обязан делиться на весь
примориал $P$ этих простых.

\begin{theorem}[\code{shifted\_primorial\_bound}]\label{thm:shifted_primorial_bound}
Если $P\mid a-\theta$ и $a\ne\theta$, то $P\le|a-\theta|$.
\end{theorem}

\emph{Идея доказательства.} Из $a\ne\theta$ следует $a-\theta\ne0$, значит $|a-\theta|>0$, и
\code{Int.le\_of\_dvd} даёт $P\le|a-\theta|$. \green{} Это \emph{закон дефекта}: как только примориал
перерастает активный делитель, нетривиального беспошлинного прохода нет
(\code{late\_boundary\_not\_free}). Он точно картирует стену чётности, но не обходит её --- суммарная
потеря ёмкости $\prod_{q\le A}\tfrac{q-2}{q-1}\sim 1/\ln A$ (Мертенс) есть стена, которую не пробить
никакой свободной от распределения бухгалтерией. Единственный скалярный вход этого маршрута --- всё
то же чётностно-остаточное препятствие.

\subsection{SNOL: редукция к соседу ранга 1}

Вместо того чтобы победить счётный бюджет, мы редуцируем всю программу к одному алгебраическому
узлу, который \emph{не является счётным по построению}. Дефекты произведённого состояния устроены по
рангу, а спуск по рангу сводит всякий дефект к рангу~$1$, где одна сторона клина есть само активное
простое $a$, так что противоположная сторона --- сдвинутый сосед $a-2\varepsilon$. Алгебра ранга~1
(\code{rank1\_opposite}, \code{neighbour\_saturation}, \code{neighbour\_corridor\_bound}) доказана
без единого распределительного допущения. Подача того же блочного моста даёт условную теорему.

\begin{theorem}[\code{twin\_primes\_of\_SNOL}]\label{thm:twin_primes_of_SNOL}
Если выполнено \code{SNOLInput} --- при каждом $N$ носитель выше $N$, чьи bad-центры (носители
терминального сдвинуто-соседского препятствия $p\mid a-2\varepsilon$, $p\le A$) строго
малочисленнее, а всякий выживший --- близнец, --- то \code{TwinLowers.Infinite}.
\end{theorem}

\emph{Идея доказательства.} Прямой мост \cref{thm:twin_prime_conjecture_of_blocks} со входом,
поданным как SNOL, а не как four-corner; \code{finite\_contradicts\_SNOL} завершает контрапозицию.
\green{} условно относительно \red{} узла \code{SNOLInput}. Соседский аудит показывает, что у
$62$--$78\%$ активных простых сосед уже пойман малым простым, и доля \emph{растёт} с $A$: SNOL нельзя
ни доказать, ни опровергнуть счётом/ёмкостью КТО в обе стороны. Значит, он обязан опираться не на
распределение, а на \emph{евклидову родословную} $a$ --- на то, что $a$ пришёл из спуска центра
клина.

\subsection{Улов --- шаг спуска: old-peel замыкается на двигатель}

Терминал $p\mid a-2\varepsilon$ --- \emph{не} терминал. Это буквально делимость противоположной
стороны клина (\code{catch\_is\_opposite}), и она разворачивается в \emph{old-peel}: разложение
$6n-\varepsilon=p\,(6t+\delta)$, порождающее строго меньший центр $t$, с жёстким законом знака
$\delta=-\pi\varepsilon$ (\code{old\_peel\_sign}).

\begin{theorem}[\code{old\_peel\_height\_drop}]\label{thm:old_peel_height_drop}
Если $6n-\varepsilon=p\,(6t+\delta)$ при $p\ge5$, $\varepsilon,\delta\in\{\pm1\}$, $t\ge1$, $n\ge2$,
то $t<n$.
\end{theorem}

\emph{Идея доказательства.} Из $6t+\delta=(6n-\varepsilon)/p\le(6n+1)/5$ и $n\ge2$ получаем $t<n$
(\code{nlinarith}). \green{} Численно спуск ещё резче --- $t<n/5$ на $3000/3000$ реальных уловах
ранга~1. Всякий улов --- шаг вниз; поток не может съезжать вниз вечно.

\begin{theorem}[\code{no\_infinite\_old\_peel}]\label{thm:no_infinite_old_peel}
Для любой $z:\N\to\N$ со свойством \code{StrictAnti}\,$z$ --- \code{False}.
\end{theorem}

\emph{Идея доказательства.} Это \code{no\_infinite\_engine\_descent}, применённое к цепи высот
old-peel: бесконечно строго убывающей цепи в $\N$ не существует. \green{} Финальный узел SNOL тем
самым замыкается \emph{не на новую теорему о распределении, а на уже доказанный двигатель}. Остаётся
структурная посылка: несточный центр всегда регенерирует корректного преемника.

\subsection{NOPSL: поток достигает близнеца}

Абстрагируем ровно ту структуру, которую поставляет old-peel.

\begin{definition}[\code{OldPeelLedger}]\label{def:oldpeelledger}
Old-peel-гроссбух над состояниями $\sigma$ --- это кортеж $(h,\mathrm{Sink},\mathrm{Step},
\code{step\_drops},\code{regenerate})$ с $h:\sigma\to\N$, предикатами
$\mathrm{Sink},\mathrm{Step}$, законом old-peel
$\code{step\_drops}:\mathrm{Step}\,s\,s'\Rightarrow h(s')<h(s)$ и регенерацией
$\code{regenerate}:\neg\,\mathrm{Sink}(s)\Rightarrow\exists s',\ \mathrm{Step}\,s\,s'$.
\end{definition}

Структура ничего не знает о простых и делимости. Вся арифметика истрачена на два законных поля;
дальнейшее рассуждение --- чистое «высота падает» и «поток не зависает».

\begin{theorem}[\code{twin\_primes\_of\_nopsl}]\label{thm:twin_primes_of_nopsl}
Дано семейство гроссбухов $L(N)$ со стартами и отображением центра плюс условие
$\code{sink\_is\_twin}:\ L(N).\mathrm{Sink}(s)\Rightarrow N<\mathrm{center}(N,s)\wedge
\mathrm{IsTwinCenter}(\mathrm{center}(N,s))$. Тогда \code{TwinLowers.Infinite}.
\end{theorem}

\emph{Идея доказательства.} Через \code{no\_old\_peel\_sink}: в предположении отсутствия стока
\code{regenerate} плюс выбор строят бесконечную траекторию, высоты которой строго убывают
(\code{step\_drops}), вопреки \cref{thm:no_infinite_old_peel}. Значит сток есть при каждом масштабе;
\code{sink\_is\_twin} делает его центром близнецов выше $N$, а
\cref{thm:infinite_of_unbounded_centers} завершает. \green{} Полная импликация
$\code{OldPeelLedger}\Rightarrow\code{TwinLowers.Infinite}$ машинно проверена на чистой логике
строгого спуска. Из двух законных полей \code{step\_drops} --- доказанная арифметика; единственный
оставшийся вход --- структурное поле \code{regenerate}: «потоку некуда деться, кроме как вниз или в
близнеца».

\subsection{Регенерация элементарна в трёх из четырёх классов}

За словом «регенерирует» скрыта безусловная дихотомия любого центра $t$.

\begin{theorem}[\code{regeneration\_dichotomy}]\label{thm:regeneration_dichotomy}
Для всякого центра $t$ и порога $A$ верно ровно одно:
\[
\begin{aligned}
  &\mathrm{Twin}(t)\ \lor\ \bigl(\neg\mathrm{OldFree}(A,t)\wedge\exists q\ \text{простое},\,5\le q\le A,\,q\mid 6t\pm1\bigr)\\
  &\lor\ \bigl(\mathrm{OldFree}(A,t)\wedge(\neg(6t-1)\ \text{простое}\ \lor\ \neg(6t+1)\ \text{простое})\bigr).
\end{aligned}
\]
\end{theorem}

\emph{Идея доказательства.} \code{by\_cases} по $\mathrm{OldFree}\,A\,t$: не-old-free даёт делитель
old-peel (\code{not\_oldfree\_gives\_peel}); old-free расщепляется по простоте, а составная old-free
сторона несёт делитель $>A$ (\code{composite\_oldfree\_has\_big\_divisor}), ребро активного спуска.
\green{} Полна и безусловна. В классах (2),(3) предъявлено исходящее ребро спуска; в (1) --- законный
сток-близнец. Нетривиальный остаток --- класс (4), fan-in/Hall --- утверждение о \emph{множестве}
родословных, не выводимое из $\mathrm{OldFree}$.

Чётностно-плотностные входы, некогда выглядевшие распределительными, погашаются
\emph{экзистенциально}: чистый старт выше любого $N$ выписывается формулой $m=(N+1)P_A$
(\code{carrier\_nonempty\_above}), а чистый сток ниже $A^2$ автоматически есть близнец.

\begin{theorem}[\code{sink\_is\_twin}]\label{thm:sink_is_twin}
Если обе стороны удовлетворяют $2\le 6m\pm1<A^2$ и ни одно простое $q\le A$ не делит ни одну из них,
то $\mathrm{TwinCenterZ}\,m$.
\end{theorem}

\emph{Идея доказательства.} Число без простых делителей $\le A$ и ниже $A^2$ не может быть составным
(его наименьший простой делитель удовлетворял бы $p^2\le n<A^2$, значит $p<A$). Две копии
\code{oldfree\_below\_sq\_prime}, по одной на сторону. \green{} Чистота исключает малые делители,
порог $A^2$ --- большие; для составной стороны места не остаётся, и никакая плотность простых не
привлекается.

\subsection{Чистый граф: закрытие бреши «сток = чистый»}

Спуск из чистого старта может добраться до малого нечистого близнеца ($18\mapsto(107,109)$), к
которому лемма привязки неприменима. Починка --- по определению: сток определён лишь внутри чистого
графа, а ребро в нечистую цель --- \emph{выход на границу}, отправляемый в гроссбух дефектов.

\begin{theorem}[\code{clean\_center\_outcome}]\label{thm:clean_center_outcome}
Для всякого чистого центра $m$ верно ровно одно:
\[
  \mathrm{TwinCenterZ}\,m\ \lor\ (\exists n,\ \mathrm{CleanActiveEdge}\,A\,m\,n)\ \lor\
  (\exists n,\ \mathrm{BoundaryExit}\,A\,m\,n)\ \lor\ \mathrm{CleanSink}\,A\,m.
\]
\end{theorem}

\emph{Идея доказательства.} Вложенный разбор случаев по близнецовости и по существованию активного
ребра, закрываемый \code{active\_edge\_clean\_or\_boundary} (исключённое третье на
$\mathrm{Clean}\,A\,n$). \green{} Классификация корректно исключает нечистого малого близнеца из
случая стока: $18$ оседает в граничной ветви, не в стоке. Но граф далёк от замкнутости --- у $59\%$
чистых центров все активные рёбра ведут на границу. Весь остаток концентрируется в одном \red{}
структурном входе: граничное состояние обязано регенерировать.

\subsection{Глобальный узел и скелет Дирихле}

Поточечная граничная регенерация \emph{ложна} ($18\mapsto(107,109)$ не имеет преемника, но и не
является корректным стоком). Поточечно доказуема лишь дихотомия исхода
(\code{boundary\_exit\_decomposes}): нечистая цель либо продолжает спуск old-peel, либо оседает в
конечный старый абсорбер ниже порога $M_0$. Нагрузка тогда переходит на всё множество свежих
стартов. Если новых близнецов выше $M_0$ нет, бесконечно много чистых стартов обязаны осесть в
конечное число старых абсорберов --- fan-in, который числа фиксируют вплоть до $570\to1$.

\begin{theorem}[\code{global\_absorber\_forces\_engine}]\label{thm:global_absorber_forces_engine}
Пусть $S$ бесконечно, $\mathrm{Codom}$ конечно, а $\mathrm{key}:\alpha\to\mathrm{Codom}$ --- паспорт
старта. Дан насос
\begin{equation}\label{eq:pump}
  \forall\gamma_1\gamma_2,\ \gamma_1\ne\gamma_2\to\mathrm{key}\,\gamma_1=\mathrm{key}\,\gamma_2\to\code{Engine}.
\end{equation}
Тогда \code{Engine}.
\end{theorem}

\emph{Идея доказательства.} Честный Дирихле: без двигателя \code{pump} запрещён, значит \code{key}
инъективно на $S$, значит бесконечное множество инъективно вкладывается в конечный тип
(\code{Set.Finite.of\_finite\_image}) --- невозможно. Коллизия $\gamma_1\ne\gamma_2$ с равным
паспортом --- \green{} доказанная часть; вся оставшаяся нагрузка в \red{} гипотезе \cref{eq:pump},
узле Холла «две различные родословные с одним паспортом $\Rightarrow$ двигатель». В сборке с
доказанной невозможностью двигателя \code{no\_global\_absorption\_under\_epmi} вынуждает новый
близнец выше $M_0$ --- условно относительно узла.

\begin{remark}
Между главами граничный узел штурмовался со многих сторон --- взвешенная долговая энергия, башни
обратного предела, барьеры прыжка/среза, оплаченная динамика, гигиена замкнутой вселенной и прочтение
$\mathrm{P}$ против $\mathrm{NP}$. Каждая обёртка машинно аудировалась; всякая либо оказывалась
циклической (редукция становилась эквивалентной цели), либо пустой (её посылка недостижима), либо
дезориентированной (вся арифметическая динамика идёт вниз, тогда как продвижение --- вверх). Ни одна
обёртка не опустила содержание ниже узла, и два внутренних аудита поймали настоящие эпизоды пустоты
прежде, чем те дошли до витрины. Красная линия удержана: единственная нестандартная аксиома ---
\axiom{}, а единственный \code{sorry} --- \code{twin\_prime\_conjecture}.
\end{remark}

\subsection{Жёсткое замыкание: спуск без цикла}

Старый довод «нет стоков $\Rightarrow$ ориентированный цикл $\Rightarrow$ двигатель $\Rightarrow\bot$»
логически избыточен. Над \code{HeightGraph} (высота, предикат близнеца, шаг с
$\code{step\_drops}:\mathrm{Step}\,s\,t\Rightarrow h(t)<h(s)$) строгое падение высоты исключает
бесконечную цепь напрямую.

\begin{theorem}[\code{reaches\_twin}]\label{thm:reaches_twin}
Если $G$ регенерирует (у всякого несблизнецового состояния есть убывающий преемник), то
$\forall s,\ \exists t,\ \mathrm{Twin}\,t$.
\end{theorem}

\emph{Идея доказательства.} Сильная индукция по высоте старта: несблизнецовое состояние шагает вниз к
строго меньшей высоте, а спускаться вечно в $\N$ нельзя. \green{} Циклическая ветвь и её жёсткая
сигнатура исчезают как отдельные входы; \code{no\_cycle} хранит невозможность цикла лишь как
переупаковку фундированности. Остаётся построить абстрактное ребро арифметикой.

\begin{theorem}[\code{cofactor\_is\_center}]\label{thm:cofactor_is_center}
Для $t\ge1$, $\eta\in\{\pm1\}$, простого $q\ge5$, взаимно простого с $6$, при $q\mid 6t+\eta$
существуют $t',\eta'$ с $\eta'\in\{\pm1\}$, $6t'+\eta'=(6t+\eta)/q$ и $0\le t'<t$.
\end{theorem}

\emph{Идея доказательства.} Кофактор $c=(6t+\eta)/q$ вынужденно взаимно прост с $6$ (таковы и
$6t+\eta$, и $q$, так что $c\bmod6\in\{1,5\}$), положителен и $<t$, ибо $5c\le 6t+1$. \green{}
Проверено на $100\%$ на $307010$ случаях. Делимость тем самым автоматически превращается в
корректный меньший центр --- ребро \emph{построено}, а не постулировано. Единственный оставшийся вход
--- полнота дихотомии на масштабе носителя (\code{regenerates\_needs\_target\_center}): что fan-in/Hall
не вводит неклассифицируемого терминала.

\subsection{Машина product-core: спуск по рангу \texorpdfstring{$4\to1$}{4->1}}

Разделяющий масштаб $6X_A+1<P_A$ (примориал обгоняет носителевую границу, достижимо с экспоненциальным
запасом уже при $A\approx50$) делает грубый паспорт $a\bmod P_A$ инъективным на легальном слое.
Перенесённое в индукцию по рангу, это закрывает \code{no\_productHall} на бесхвостом
\code{RankNode} $r$ --- знак плюс $\mathrm{factors}:\mathrm{Fin}\,r\to\N$ --- где экстенсиональность
есть подлинная теорема (\code{no\_mismatch\_core\_eq}), а не гипотеза.

\begin{theorem}[\code{no\_productHall}]\label{thm:no_productHall}
Если в роли $k$ оба множителя $<P_A$ и сравнимы по модулю $P_A$, но различны, то \code{False}.
\end{theorem}

\emph{Идея доказательства.} Для $a<P_A$ имеем $\code{ZMod.val}(a\bmod P_A)=a$, так что равные вычеты
вынуждают равные множители, вопреки различию. \green{} Удаление хвоста из объекта индукции убирает
fan-in $570\to1$, а с ним и все три аудиторских дефекта (экстенсиональность, оценку остатка
\code{ambientLegal\_delete}, конечность сигнатуры при фиксированном $A$) разом.

\begin{theorem}[\code{core\_step\_proved}]\label{thm:core_step_proved}
При разделяющем масштабе $\forall r,\ 2\le r\Rightarrow\code{CoreCollision}\,X_A\,P_A\,r\Rightarrow
\code{CoreCollision}\,X_A\,P_A\,(r-1)$.
\end{theorem}

\emph{Идея доказательства.} Из равных сигнатур и различных узлов есть несовпадающая роль $k$; удалим
её. Либо остатки различны (коллизия меньшего ранга), либо совпадают, концентрируя рассогласование в
$k$ --- ProductHall, невозможный по \cref{thm:no_productHall}. \green{} База ранга~$1$ закрыта той же
арифметикой разделяющего масштаба (\code{rank\_one\_coreCollision\_absurd}), без внешнего SNOL.
Сильная индукция по рангу даёт коллизия $\Rightarrow$ \code{Engine}, а Дирихле по конечной сигнатуре
(\code{coreCollision\_of\_infinite}) изготавливает коллизию из бесконечного инъективного семейства
стартов.

\begin{theorem}[\code{product\_core\_engine\_of\_carrier}]\label{thm:product_core_engine_of_carrier}
При разделяющем масштабе, $1\le r\le4$ и бесконечном $S$ амбиент-легальных \code{RankNode} с
$\mathrm{node}$, инъективным на $S$, получаем \code{Engine}.
\end{theorem}

\emph{Идея доказательства.} Собрать экстенсиональность, оценку, конечность, спуск
(\cref{thm:core_step_proved}), арифметическую базу и Дирихле
(\code{coreCollision\_of\_infinite}). \green{} \emph{Условная} теорема: всё между носителевой
гипотезой и \code{Engine} машинно проверено. Дыра сжата в один вход, \code{CarrierData} ---
бесконечное инъективное семейство амбиент-легальных стартов фиксированного ранга.

\subsection{Факторизация и последнее звено}

Носителевая гипотеза поставляется, узел за узлом, чистой факторизационной арифметикой.

\begin{theorem}[\code{mkNode\_of\_composite}]\label{thm:mkNode_of_composite}
Пусть $1\le A$, $1<N$, $N$ составное, $N\le 6X_A+1$, $6X_A+1<A^5$, и всякий простой делитель $N$
больше $A$. Тогда для любого знака есть ранг $r$ с $2\le r\le4$ и $X:\code{RankNode}\,r$ с
$\code{AmbientLegal}\,X_A\,X.\mathrm{factors}$ и $\prod(\mathrm{ofFn}\,X.\mathrm{factors})=N$.
\end{theorem}

\emph{Идея доказательства.} Берём $L=\code{primeFactorsList}\,N$: составность даёт $|L|\ge2$
(\code{composite\_rank\_ge\_two}); чистота делает всякий множитель $>A$; а
$(A+1)^{|L|}\le\prod L=N\le 6X_A+1<A^5<(A+1)^5$ вынуждает $|L|\le4$ (\code{factor\_rank\_le\_four}).
Свидетель $N_0=N$ удостоверяет \code{AmbientLegal}. \green{} Магическое число $4$ --- арифметическое
следствие масштаба решета, а не соглашение. Соединяя эту фабрику узлов с доказанной бесконечностью
носителя:

\begin{theorem}[\code{cleanCenters\_infinite}]\label{thm:cleanCenters_infinite}
Для всякого $A$ множество $\mathrm{CleanCenters}(A)=\{m\mid\mathrm{CleanZ}\,A\,m\}$ бесконечно.
\end{theorem}

\emph{Идея доказательства.} Через \code{Set.infinite\_of\_not\_bddAbove}: при любом $N$
\code{carrier\_nonempty\_above} даёт чистый $m>N$ (кратное старого примориала). \green{} Плотность не
нужна.

\begin{theorem}[\code{engine\_of\_carrier\_and\_factorize}]\label{thm:engine_of_carrier_and_factorize}
При разделяющем масштабе, бесконечном носителе $C$ и отображениях ранга \code{rankOf} и узла
\code{mkNode}, инъективных и \code{AmbientLegal} на каждом ранге, выполнено \code{Engine}.
\end{theorem}

\emph{Идея доказательства.} Композиция \code{factorizationData\_of\_carrier} (бесконечный слой одного
ранга, выбранный бесконечным Дирихле \code{exists\_infinite\_fiber}) с
\cref{thm:product_core_engine_of_carrier}. \green{} Бесконечность носителя, выбор слоя и вся насосная
машина доказаны полностью. Единственный вход --- факторизационные отображения.

\subsection{Единственный открытый узел}

Здесь нужна точность, иначе редукция выдаётся за доказательство. Легальность
$\code{AmbientLegal}\,X_A$ требует, чтобы множители делили общую верхнюю сторону $N_0\le6X_A+1$, то
есть $m\lesssim X_A$: \emph{конечное} горизонтальное окно. Но \cref{thm:cleanCenters_infinite} даёт
бесконечность \emph{по вертикали}, $m\to\infty$, неизбежно выше $X_A$, где сертификат --- а с ним и
оценка ранга $\le4$ --- отказывает. Два доказанных конца не сращиваются при фиксированном масштабе.
Требуется механизм, поставляющий бесконечно много легальных узлов на \emph{растущем} масштабе, где
$A$ и $X_A$ растут вместе с $m$.

\begin{conjecture}[\code{GlobalOldAbsorption}, единственный открытый узел]\label{conj:globaloldabsorption}
\red{} Существует бесконечное семейство чистых центров $m$, когерентно при $m\to\infty$, каждый с
составной стороной, допускающей легальную факторизацию фиксированного ранга $1\le r\le4$ в простые
$>A(m)$ при масштабе $A(m),X_A(m)$, сохраняющем разделяющий масштаб $6X_A(m)+1<P_{A(m)}$, причём
различные центры дают различные узлы.
\end{conjecture}

Это последняя обязанность Step00. Она \emph{неснижаема} --- по силе равна самой гипотезе о близнецах
--- и не может быть закрыта изнутри системы: закрыть её --- значит запустить тот самый двигатель,
невозможность которого лежит в основании. Барьер структурен (горизонтальное окно обрезает любой
фиксированный масштаб), а не артефакт одного маршрута: счёт, спуск EPMI и four-corner все сходятся
сюда, и всякая попытка-обёртка машинно аудировалась обратно к нему. Соответственно машина фиксирует
ровно одну заражённую \code{sorry} декларацию, \code{twin\_prime\_conjecture}, замыкающуюся через
\cref{thm:engine_of_carrier_and_factorize} на этот узел.

\begin{remark}
Повторим честный вердикт. Цепь
\[
\begin{aligned}
  &\code{survivor\_of\_not\_covered}\to\code{twin\_prime\_conjecture\_of\_blocks}\\
  &\to\code{twin\_primes\_of\_nopsl}\to\code{product\_core\_engine\_of\_carrier}\\
  &\to\code{engine\_of\_carrier\_and\_factorize}
\end{aligned}
\]
\green{} машинно проверена и свободна от распределения от края до края. Гипотеза о близнецах
\emph{не} доказана. Доказана \emph{редукция} её к единственному узлу
\cref{conj:globaloldabsorption}, помеченному \red{}, который \cref{sec:firstcause} принимает как
первопричину \axiom{} извне системы.
\end{remark}
